\section{Totality of circular proofs}
\label{sec:totality}



\comment{
Minor things:
\begin{itemize}
    \item explain difference with interpretation of fixed points in models of F
    \item mention, maybe in previous section, that untyped terms are Turing-complete.
    \item need to mention that similar semantic techniques in literature, e.g. Mendler, Barendregt... this also highlights differences with F.
    \item make a comment about (weak) extensionality: we do not take full extensionality to keep equational theory semi-recursive, avoiding extensionality-elimination techniques à la Luckhardt.
    \item notation: $\clomuLJneg$ should have a default equational theory, maybe put $\eta$ in sub/superscript
\end{itemize}
}


In this section we provide a semantics for (circular) proofs, using computability theoretic tools.
Our aim is to show that $\cmuLJ$ represents only total functions on $\Nat$ (Corollary~\ref{cor:cmuLJ-represents-only-total-fns}), by carefully extending circular proof theoretic techniques to our semantics.

Throughout this section we shall only consider types formed from $\N, \times, \arrow, \mu$, unless otherwise indicated.

\subsection{A type structure of regular coterms}
% We shall now turn to building up a (impredicative) higher-order computability model in order to obtain our totality proof for $\cmuLJneg$ (and hence $\cmuLJ$, for functions on natural numbers).
We shall define a type structure whose domain will be contained within $\closure \cmuLJneg$. 
Before that, it will be convenient to have access to a notion of a `good' set of terms.
% 

A \emph{(totality) candidate} is some $A\subseteq \closure \cmuLJneg$ that is closed under $\convertsnoreceta$.
We henceforth expand our language of types by including each candidate $A$ as a type constant. 
% the clause:
% \begin{itemize}
% \item each candidate $A$ is a type constant.
% \end{itemize}
% (Types are still closed under $\times, \arrow, \mu$).
% 
An immediate albeit powerful observation is that the class of candidates forms a complete lattice under set inclusion.
This justifies the following definition of our type structure:
\begin{definition}
    [Type structure]
    \label{def:hr}
    For each type $\sigma$ we define $\hr\sigma \subseteq \closure \cmuLJneg$ by:
    \[
    \begin{array}{r@{\ := \ }l}
         \hr A & A \\
         \hr \N & \{ t  \ \vert \ \exists n \in \Nat .\, t\convertsnoreceta \numeral n\} \\
         \hr {\sigma \times \tau} & \{ t \ \vert\ \proj 0 t \in \hr \sigma \text{ and }  \proj 1 t \in \hr \tau\} \\
         \hr {\sigma \to \tau} & \{ t \ \vert\ \forall s \in \hr \sigma. \, ts \in \hr\tau\} \\
         \hr {\mu X\sigma(X)} & \bigcap \{ A \text{ a candidate} \ \vert\ \hr{\sigma(A)}\subseteq A\}
    \end{array}
    \]
    We write $\hr{\sigma(\cdot)}$ for the function on candidates $A\mapsto \hr{\sigma(A)}$.
\end{definition}
As we shall see, the interpretation of $\mu$-types above is indeed a least fixed point of the corresponding operation on candidates.

\begin{remark}
    [Alternative SO interpretation]
    Recalling the second-order interpretation of $\mu$-types, \Cref{rem:mulj-frag-F}, an alternative definition of $\hr {\mu X \sigma(X)}$ could be $\bigcap_A (\sigma(A)\to A) \to A$.
    This gives rise to a different type structure, indeed similar to the realisability model we give later in \Cref{sec:realisability}.
    However such a choice does not allow us to readily interpret $\clocmuLJneg$: the totality argument in this section, \Cref{thm:prog-implies-hr}, crucially exploits the fact that $\mu$-types are interpreted as bona-fide least fixed points.
\end{remark}

A routine but important property is:
\begin{proposition}
    [Closure under conversion]
    \label{prop:hr-closed-under-conversion}
    If $t \in \hr \tau$ and $t \convertsnoreceta t'$ then $t'\in \hr\tau$.
\end{proposition}
% {\begin{proof}
% It suffices to show that the reduction rules in~\Cref{fig:reduction-ljneg},~\Cref{fig:reduction-mu-norec},~\Cref{fig:native-n-rules-mulj}, and the rule ($\extensional$) preserve the type structure. We just show the case of the equations for the fixed point combinators. This can be easily established by proving the following properties:
% \[
% \begin{array}{rcl}
%  \rr \mu & \dfn& \hr{(\vec \tau \to \sigma(\mu X. \sigma))\to \vec \tau \to \mu X. \sigma}\\
%     \lr{\mu'} &\dfn&  \hr{(\vec \tau\to \sigma (\mu X. \sigma)\to \gamma) \to \vec \tau \to \mu X. \sigma \to \gamma} 
% \end{array}
% \]
% where $\vec \tau= \tau_1, \ldots, \tau_n$. 
% Concerning $\rr \mu$, suppose $t \in \hr{\vec \tau \to \sigma(\mu X. \sigma)}$ and $\vec s= s_1, \ldots, s_n$ is such that $s_i \in \hr{\tau_i}$. Then $t \, \vec s \in \hr{\sigma(\mu X. \sigma)}$. By definition, $\hr{\sigma(\mu X. \sigma)} \subseteq \hr{\mu X. \sigma}$. So, $t\, \vec s\in \hr{\mu X. \sigma}$. $\rr{\mu}\, t\, \vec s\convertsnoreceta t\, \vec s$. As for the case $\rr \mu$, suppose $t \in \hr{\vec \tau\to \sigma (\mu X. \sigma)\to \gamma}$, $\vec s= s_1, \ldots, s_n$ is such that $s_i \in \hr{\tau_i}$, and $u \in \hr{\mu X. \sigma}$. By definition, $u \in \hr{\mu X. \sigma}=\hr{\sigma(\mu X. \sigma)}$, so $ t\, \vec s\, u\in \hr{\gamma}$. But $\lr {\mu'}\, t\, \vec s\, z \convertsnoreceta t\, \vec s\, u$.
% \end{proof}
% }
\begin{proof}
    By induction on the structure of $\tau$.
    The base cases when $\tau$ is a candidate $A$ or the type $\nat$ follow immediately from the definitions. For the remaining cases:
    \begin{itemize}
        \item If $\tau = \tau_0 \times \tau_1$, then $\proj i t' \convertsnoreceta \proj i t \in \hr{\tau_i}$ by IH, for $i=0,1$, so indeed $t'\in \hr \tau$.
        \item If $\tau = \tau_0 \to \tau_1$ and $s \in \tau_0$, then $t's \convertsnoreceta ts \in \hr{\tau_1}$ by IH, so indeed $t'\in \hr\tau$.
        \item If $\tau = \mu X \tau'(X)$ and $A$ is a candidate with $\hr{\sigma(A)} \subseteq A$, then $t' \in A$ by IH, so indeed $t'\in \hr\tau$.\qedhere
    \end{itemize}
\end{proof}


Let us point out that this immediately entails, by contraposition and symmetry of $\convertsnoreceta$, closure of \emph{non-elementhood} of the type structure under conversion: if $t\notin \hr \tau$ and $t\convertsnoreceta t'$ then also $t'\notin \hr\tau$.
% This simple observation is critical for our main totality argument, which proceeds by contradiction.


% We shall discuss in more detail in the next subsection.
The main result of this section is:
\begin{theorem}
[Interpretation]
\label{thm:prog-implies-hr}
For any $\cmuLJneg$-coderivation $\der:\Sigma \seqar \tau$ and $\vec s \in \hr\Sigma$ we have $\der\vec s \, \in  \hr{ \tau}$.
\end{theorem}

\anupam{I moved the proof of this to the end of section, as it is way too early here. We need the notions built up over the next subsections for it to make sense. Added sentence below.}

The rest of this section is devoted to proving this result, but before that let us state our desired consequence:

% Note that, by simply unfolding the definition we obtain:

% \begin{observation}
%     [Totality at type 1]
%     If $t \in \hr{\vec \N \to \N}$ then $ \{(\vec m,n) \in \vec \Nat \times \Nat \ \vert \ t \numeral{\vec m} \, \convertsnorec \numeral n\} $ is the graph of a (total) function
%     $\vec \Nat \to \Nat$.\footnote{The lengths of vectors match in the expected way.} 
% \end{observation}

% Thus, from the Interpretation Theorem we immediately obtain:

\begin{corollary}
\label{cor:cmuLJ-represents-only-total-fns}
$\cmuLJneg$ represents only total functions on $\Nat$ with respect to $\convertsnoreceta$.
\end{corollary}
\begin{proof}
[Proof idea]
Consider a $\cmuLJneg$-coderivation $P: \vec \N \seqar \N$.
    By \Cref{thm:prog-implies-hr} and closure under conversion, Proposition~\ref{prop:hr-closed-under-conversion} (namely applying $\cut$-reductions), we have for all $\vec m \in \Nat$ there is $n\in \Nat$ with  $P\numeral{\vec m} \convertsnoreceta \numeral n$.
\end{proof}


\subsection{Montonicity and transfinite types}
\label{subsec:monotonicity-and-transfinite-types}
To prove our main Interpretation Theorem, we shall need to appeal to a lot of background theory on fixed point theorems, ordinals and approximants, fixed point formulas, and cyclic proof theory. 
In fact we will go on to formalise this argument within fragments of second-order arithmetic.
% Since we will actually formalise these intermediate results in fragments of second-order arithmetic in the next section, the exposition herein will be rather brief.



Since the class of candidates forms a complete lattice under set inclusion, we can specialise the well-known Knaster-Tarski fixed point theorem:
\begin{proposition}
    [Knaster-Tarski for candidates]
    \label{prop:kt-on-candidates}
    Let $F$ be a monotone operation on candidates, with respect to $\subseteq$.
    $F$ has a least fixed point
     $\mu F = \bigcap \{A \text{ a candidate}\ \vert\ F(A)\subseteq A\}$.
\end{proposition}

At this point it is pertinent to observe that the positivity constraint we impose for fixed point types indeed corresponds to monotonicity of the induced operation on candidates with respect to our type structure:
\begin{lemma}
[Monotonicity]
\label{lem:mon-pos-types-in-hr}
Let $A\subseteq B$ be candidates.
\begin{enumerate}
    \item 
If $\sigma(X)$ is positive in $X$ then $\hr{\sigma(A)}\subseteq \hr{\sigma(B)}$; 
\item If $\sigma(X)$ is negative in $X$ then $\hr{\sigma(B)}\subseteq \hr{\sigma(A)}$.
\end{enumerate}
\end{lemma}
These properties are proved (simultaneously) by a straightforward induction on the structure of $\sigma(X)$. \todo{give proof?}
Note that, by the Knaster-Tarski fixed point theorem this yields:
\begin{proposition}
[``Fixed points'' are fixed points]
\label{prop:hr-mu-is-lfp}
    $\hr{\mu X \sigma(X)}$ is the least fixed point of $\hr{\sigma(\cdot)}$ on candidates.
\end{proposition}

We will need to appeal to an alternative characterisation of fixed points via an inflationary construction, yielding a notion of `approximant' that:
\begin{enumerate}[(a)]
    \item allows us to prove the Interpretation Theorem by reduction to well-foundedness of approximants (or, rather, the ordinals that index them); and 
    \item allows a logically simpler formalisation within second-order arithmetic, cf.~Section~\ref{sec:rev-math-kt}, crucial for obtaining a tight bound on representable functions,
    % in Section~\ref{sec:formalised-totality}.
\end{enumerate}

\begin{definition}
    [Approximants]
    Let $F$ be a monotone operation on candidates, with respect to $\subseteq$.
    For ordinals $\alpha$ we define $F^\alpha(A)$ by transfinite induction on $\alpha$:
    \begin{itemize}
        \item $F^0 (A) \dfn \emptyset$
        \item $F^{\succ \alpha}(A) \dfn F(F^\alpha(A))$
        \item $F^{\lambda}(A) \dfn \bigcup\limits_{\alpha<\lambda} F^\alpha(A)$, when $\lambda$ is a limit ordinal.
    \end{itemize}
\end{definition}
For our purposes we will only need the special case of the definition above when $A=\emptyset$.
Writing $\Ord$ for the class of all ordinals, the following is well known:
\begin{proposition}
    [Fixed points via approximants]
    \label{prop:fixed-points-via-approximants}
Let $\mu F$ be the least fixed point of a monotone operation $F$. 
% be as in Proposition~\ref{prop:kt-on-candidates}.
    We have
    $\mu F = \bigcup\limits_{\alpha \in \Ord} F^\alpha(\emptyset)$.
\end{proposition}
    
From here it is convenient to admit formal type expressions representing approximants.
\begin{convention}
    [Transfinite types]
    We henceforth expand the language of types to be closed under:
\begin{itemize}
    \item for $\sigma(X)$ positive in $X$, $\alpha$ an ordinal, $\sigma^\alpha(\tau)$ is a type.
\end{itemize}
Again we shall only need the special case of $\tau = \emptyset$ for our purposes.
We call such types \emph{transfinite} when we need to distinguish them from ordinal-free types.
\end{convention}

\begin{definition}
    [Type structure, continued]
    We expand Definition~\ref{def:hr} to account for transfinite types by setting $\hr{\sigma^\alpha(\tau)} \dfn \hr{\sigma(\cdot)}^\alpha (\hr \tau)$. 
    In other words:
\begin{itemize}
    \item $\hr{\sigma^0(\tau)} := \hr\tau$
    \item $\hr{\sigma^{\succ \alpha}(\tau)}:= \hr{\sigma(\hr{\sigma^\alpha(\tau)})}$
    \item $\hr{\sigma^\lambda (\tau)}:= \bigcup\limits_{\alpha<\lambda} \hr{\sigma^\alpha(\tau)}$, when $\lambda$ is a limit ordinal.
\end{itemize}
\end{definition}

We have immediately from Proposition~\ref{prop:fixed-points-via-approximants}:
\begin{corollary}
\label{cor:hr-mu-as-limit-of-approximants}
$\hr{\mu X \sigma(X)} = \bigcup\limits_{\alpha \in \Ord} \hr{\sigma^\alpha(\emptyset)}$.
\end{corollary}



\subsection{Ordinal assignments}


We shall write $\sigma \subform \tau$ if $\sigma$ is a subformula of $\tau$.


The \emph{Fischer-Ladner (well) preorder}, written $\flleq$,
is the smallest extension of $\subform$, restricted to closed formulas, satisfying
$\sigma(\mu X \sigma(X)) \flleq \mu X \sigma(X)$.
We write $\sigma \fleq \tau$ if $\sigma \flleq \tau \flleq \sigma$ and $\sigma \flnleq \tau$ if $\sigma \flleq \tau \not\flleq \sigma$.
Note that $\fleq$-equivalence classes are naturally (well) partially ordered by $\flleq$.

The \emph{Fischer-Ladner closure} of a type $\tau$, written $\fl \tau$, is the set $\{\sigma \mid \sigma \flleq \tau\}$.
Note that $\fl \tau$ is the smallest set of closed types closed under subformulas and, whenever $\mu X\sigma(X) \in \fl \tau$, then also $\sigma(\mu X\sigma(X)) \in \fl \tau$; this is necessarily a finite set.

\begin{definition}
[Priority]
We say that a type $\sigma$ has higher \emph{priority} than a type $\tau$, written $\sigma>\tau$, if $\tau \flnleq \sigma$ or $\sigma \fleq \tau $ and $\sigma \subform \tau$.

\end{definition}
Note that $<$ is a (strict) well partial order on types. 
The priority order is commonly used in modal fixed point logics, e.g., \cite{Studer08:mu-calc,Doumane17thesis,venema2008lectures}.
% , as it determines dependencies between fixed point formulas at the level of their bound variables:
% \begin{remark}
%     Assuming each $\mu$ binds a distinct variable, $\mu X \sigma < \mu Y \tau \flleq \mu X \sigma$ just if $Y$ occurs free in $\sigma$ (but $X$ does not occur free in $\tau$).    
% \end{remark}

\begin{convention}
In what follows, we shall assume an arbitrary extension of $<$ to a \emph{total} well order.    
\end{convention}




Let $\tau$ be a type whose $<$-greatest fixed point subformula occurring in \emph{positive} position is $\mu X \sigma(X)$.
We write $\tau^\alpha$ for $\tau[\sigma^\alpha(\emptyset) / \mu X\sigma(X)]$, i.e.\ $\tau$ with each occurrence of $\mu X\sigma(X)$ in positive position replaced by $\sigma^\alpha(\emptyset)$.
$\tau_\alpha$ is defined the same way by for the $<$-greatest fixed point subformula occurring in \emph{negative} position.

Note that, if $\tau$ has $n$ fixed point subformulas in positive position and $\vec \alpha = \alpha_1 ,\dots, \alpha_n$ then $\tau^{\vec \alpha} = (\cdots ((\tau^{\alpha_1})^{\alpha_2})\cdots)$ is a positive-$\mu$-free (transfinite) type.
Similarly for negatively occurring fixed points.
We shall call such sequences \emph{(positive or negative) assignments} (respectively).

We shall order assignments by a lexicographical product order, i.e.\ by setting $\vec \alpha < \vec \beta$ when there is some $i$ with $\alpha_i < \beta_i$ but $\alpha_j = \beta_j$ for all $j<i$. 
Note that this renders the order type of $\vec \alpha$ simply $\alpha_n \times \cdots \times \alpha_1$, but it will be easier to explicitly work with the lexicographical order on ordinal sequences.


By the Monotonicity Lemma~\ref{lem:mon-pos-types-in-hr} we have:
\begin{proposition}
[Positive and negative approximants]
\label{prop:pos-neg-approximants}
If $t \in \hr\tau$ there are (least) ordinal(s) $\vec \alpha$ s.t.\ $t \in \hr{\tau^{\vec \alpha}}$.
% 
Dually, if $t \notin \hr \tau$ there are (least) ordinal(s) $\vec \alpha$ s.t.\ $t \notin \hr{\tau_{\vec \alpha}}$.
\end{proposition}

\anupam{commented `approximants are successors' below, now obsolete}
Since the ordinals given by the above Proposition are points of first entry for an element into a fixed point,
note that, for $\vec \alpha$ as in the Proposition above, each $\alpha_i$ must be a successor ordinal.
% \begin{remark}
% [Approximants are successors]
% \label{obs:approxs-are-successors}
% Let $\vec \alpha$ be as in the Proposition above. Then each $\alpha_i$ is a successor ordinal.
% \end{remark}



% \begin{definition}
% [Assignments]
% Let $\tau$ be a type whose positively occurring fixed points are $>$-enumerated $\mu X_1 \sigma_1(X_1) > \cdots > \mu X_n \sigma_n (X_n)$.
% A \emph{(positive) (ordinal) assignment} $\alpha$ for $\tau$ is a finite list of ordinals $\alpha = (\alpha_1, \dots , \alpha_n)$, that we shall simultaneously identify with the ordinal $\alpha_1 \times \cdots \times \alpha_n$ whenever it is convenient (so that assignments are well-ordered by ordinal comparison).
% We shall write $\tau^\alpha$ for the (transfinite) type obtained by replacing each $\mu X_i \sigma_i(X_i)$ by $\sigma_i^{\alpha_i} (\emptyset)$.
% \oldtodo{this needs to be defined a little more carefully to reflect the order of substitutions. make a definition of single substitution, then make it sequential.}


% A \emph{(negative) (ordinal) assignment} is defined the same way but for negatively occurring fixed points.
% When writing $\tau^\alpha$ we shall always make it clear in advance whether we are replacing positively or negatively occurring fixed points of $\tau$. \oldtodo{ambiguity not ideal, consider making clearer if needed}
% \end{definition}


\subsection{Reflecting non-totality in rules of $\muLJnegnorec$}

% At this juncture one natural approach to proving the Interpretation Theorem would be to directly assign ordinals to type occurrences in (co)derivations.
% However this does not seem to work as usual soundness/totality arguments for non-wellfounded proofs use an argument by contradiction to construct a `countermodel' or `non-total' branch \todo{cite arithmetic/T papers among others}, and it is only at this point that we may call on an appropriate ordinal assignment.
% For this reason we must define our non-total branch accounting for all possible ordinal assignments.
% Note that this is similar to how soundness arguments for cyclic proofs for the modal mu-calculus proceed. \todo{cites}

\anupam{commented some narrative above}


Before giving our non-total branch construction, let us first make a local definition that will facilitate our construction:
\begin{definition}
    [Reflecting non-totality]
    \label{def:refl-non-totality}
    Fix a $\muLJnegnorec$-step,
    \[
    \vliiinf{ \rrule}{}{\Sigma \seqar \tau}{\Sigma_0 \seqar \tau_0}{\cdots}{\Sigma_{n-1} \seqar \tau_{n-1}}
    \] 
    for some $n\leq 2$ and regular coderivations $\der_i:\Sigma_i\seqar \tau_i$. 
    
    For $\vec s \in \hr{\Sigma}$ s.t.\ $ \rrule {\vec \der} \vec s \notin \hr \tau$,
    we define a premiss $\Sigma' \seqar \tau'$ and $\der':\Sigma'\seqar \tau'$ and some inputs $\vec s' \in \hr{ \Sigma'}$ such that $\der'\vec s'
    \notin \hr{\tau'}$
    as follows:
\begin{itemize} 
\item $\rrule $ cannot be $\id$ since $\id s \convertsnoreceta s$.
\item If $\rrule $ is $\vlinf{\ex}{}{\Gamma, \rho, \sigma, \Delta \seqar \tau}{\Gamma , \sigma, \rho , \Delta \seqar \tau}$ and $\vec s = (\vec r, r , t,\vec t)$ with $\vec r \in \hr \Gamma, r\in \hr \rho , s \in \hr \sigma, \vec t \in \hr \Delta$, we set $(\Sigma'\seqar \tau'):= (\Gamma, \sigma, \rho, \Delta \seqar \tau)$, $\der':=\der_0$ and $\vec s':= (\vec r,t,r,\vec t)$.
\item If $\rrule $ is $\vlinf{\wk}{}{\Sigma_0,\sigma \seqar \tau}{\Sigma_0 \seqar \tau}$ and $\vec s = (\vec s_0,s) $ with $\vec s_0 \in \hr{\Sigma_0}$ and $s \in \hr \sigma$ then we set $\der':= \der_0$, $(\Sigma'\seqar \tau'):= (\Sigma_0 \seqar \tau)$ and $\vec s' := \vec s_0$. 
\item If $\rrule $ is $\vlinf{\contr}{}{\Sigma_0 , \sigma\seqar \tau}{\Sigma_0, \sigma, \sigma \seqar \tau}$ and $\vec s = (\vec s_0,s)$ with $\vec s_0 \in \hr {\Sigma_0}, s\in \hr \sigma$, then we set $(\Sigma' \seqar \tau'):= \allowbreak (\Sigma_0\seqar \tau)$, $\der':= \der_0$ and $\vec s' := (\vec s_0,s,s)$.
\item If $\rrule $ is 
\(
\vliinf{\cut}{}{\Sigma_0,\Sigma_1 \seqar \tau}{\Sigma_0 \seqar \tau_0}{\Sigma_1, \tau_0 \seqar \tau }
\)
and $\vec s_0 \in \hr{\Sigma_0}, \vec s_1 \in \hr{\Sigma_1}$ we have:
\[
\begin{array}{rll}
     & \cut \der_0 \der_1 \vec s_0 \vec s_1\, \notin \hr\tau \\
     \implies & \der_1  \vec s_1 (\der_0 \vec s_0) \, \notin \hr \tau & \text{by $\cut$ reduction}
\end{array}
\]
Now, if $\der_0\, \vec s_0 \ \notin \hr{\tau_0}$ then we set $(\Sigma'\seqar \tau') := (\Sigma_0 \seqar \tau_0)$, $\der' := \der_0$ and $\vec s' := \vec s_0$.
Otherwise $\der_0\, \vec s_0 \ \in \hr{\tau_0}$ so we set 
$(\Sigma'\seqar \tau') := (\Sigma_1 \seqar \tau_1)$, $\der' := \der_1$ and $\vec s' := (\vec s_1, \der_0\, \vec s_0)$.
\item $\rrule$ cannot be 
\(
\vlinf{\Nzero}{}{\seqar \nat}{}
\)
as $\numeral 0 \in \hr \nat$.
\item If $\rrule$ is
\(
\vlinf{\Nsucc}{}{\Sigma \seqar \nat}{\Sigma \seqar \nat}
\)
and $\vec s \in \hr{\Sigma}$,
\[
\begin{array}{rll}
     & \Nsucc \der_0  \vec s \, \notin \hr\nat \\
     \implies & \succ (\der_0 \vec s) \, \notin \hr \nat & \text{by $\Nsucc $ reduction} \\
     \implies & \der_0 \vec s \, \notin \hr \nat  & \text{by context closure of $\convertsnoreceta$}
\end{array}
\]
so we set $(\Sigma'\seqar \tau'):= (\Sigma \seqar \nat)$, $\der' := \der_0$ and $\vec s' := \vec s$.
\item If $\rrule$ is
\(
\vliinf{\Ncnd}{}{\Sigma_0,\nat \seqar \tau}{\Sigma_0 \seqar \tau}{\Sigma_0,\nat \seqar \tau}
\)
and $\vec s_0 \in \hr{\Sigma_0}$, $s \in \hr \nat$, then $s\convertsnoreceta \numeral n$ for some $n\in \Nat$ by definition of $\hr\nat$.
If $n=0$ then,
\[
\begin{array}{rll}
     &  \Ncnd \der_0 \der_1 \vec s s \, \notin \hr \tau \\
     \implies & \Ncnd \der_0\der_1 \vec s \numeral 0 \, \notin \hr \tau & \text{by closure under conversion} \\
     \implies & \der_0 \vec s \, \notin \hr \tau & \text{by $\Ncnd$ reduction}
\end{array}
\]
so we set $(\Sigma'\seqar \tau'):= (\Sigma_0 \seqar \tau)$, $\der':= \der_0$, $\vec s' := \vec s$.

Otherwise if $n= n'+1$ then,
\[
\begin{array}{rll}
     & \Ncnd \der_0 \der_1 \vec s s \, \notin \hr \tau \\
     \implies & \Ncnd \der_0  \der_1 \vec s  (\succ \numeral n') \, \notin \hr \tau & \text{by closure under conversion} \\
     \implies & \der_1  \vec s \numeral n' \, \notin \hr \tau & \text{by $\Ncnd$ reduction}
\end{array}
\]
so we set $(\Sigma'\seqar \tau'):= (\Sigma_0, \nat \seqar \tau)$, $\der':= \der_1$, $\vec s' := (\vec s,\numeral n')$.
        \item If $\rrule$ is
        \(
        \vliinf{\times r }{}{\Sigma_0, \Sigma_1 \seqar \tau_0 \times \tau_1}{\Sigma_0 \seqar \tau_0}{\Sigma_1 \seqar \tau_1}
        \)
        and $\vec s_0 \in \hr{\Sigma_0}, \vec s_1 \in \hr{\Sigma_1}$,
        \[
        \begin{array}{rll}
        & {\rr \times} {\der_0}{\der_1}\vec s_0 \vec s_1 \notin \hr{\tau_0\times \tau_1} & 
        % \text{by assumption}
        \\ 
        \implies & \pair{\der_0\vec s_0}{\der_1\vec s_1} \notin\hr{\tau_0\times \tau_1} & \text{by $\rr\times$ reduction} \\
        \implies & \proj i \pair{\der_0\vec s_0}{\der_1\vec s_1}  \notin \hr{\tau_i} & \text{by $\hr{\times}$} \\
        \implies & {\der_i}\vec s_i \notin \hr{\tau_i} & \text{by {$\proj i \pair{\cdot}{\cdot}$ reduction}}
        \end{array}
        \]
        for some $i\in \{0,1\}$,
        so we set $(\Sigma' \seqar \tau') := (\Sigma_i \seqar \tau_i)$, $\der' := \der_i$, $\vec s' := \vec s_i$.
    \item If $\rrule$ is $\vlinf{  {\lr \times}}{}{\Gamma, \sigma_0 \times \sigma_1 \seqar \tau  }{\Gamma, \sigma_0,\sigma_1 \seqar \tau}$ and $\vec r \in \hr{\Gamma}$ and $s \in \hr{\sigma_0 \times \sigma_1}$:
    \[
    \begin{array}{rll}
        & {\lr\times }\,  {\der_0}\,  \vec r \, s\ \notin\  \hr\tau & 
        % \text{by assumption} 
        \\
        \implies & {\der_0}\, \vec r \, \proj 0 s\, \proj 1 s \ \notin\  \hr\tau & \text{by $\lr \times$ reduction}
    \end{array}
    \]
    Now, by definition of $\hr{\sigma_0 \times \sigma_1}$ we have indeed $\proj i s \in \hr{\sigma_i}$ for $i\in \{0,1\}$,
    so we set $(\Sigma' \seqar \tau') = (\Gamma,\sigma_0,\sigma_1 \seqar \tau)$, $\der' = \der_0$ and $\vec s' = (\vec r, \proj i s)$.
\item If $\rrule$ is
    \(
    \vlinf{\rr\arrow}{}{\Sigma \seqar \rho \to \sigma }{\Sigma, \rho \seqar \sigma}
    \)
    and $\vec s \in \hr{\Sigma}$,
    \[
    \begin{array}{rll}
         & {\rr\arrow } {\der_0}\vec s\  \notin\  \hr{\rho\to \sigma} & \\
         \implies & {\rr \arrow} {\der_0} \vec s \  \notin\  \hr{(\rho \to \sigma)_{\vec \alpha}} & \text{for some least $\vec \alpha$}\\
         \implies & \rr \arrow  \der_0  \vec s\  \notin \ \hr{\rho^{\vec \alpha_0} \to \sigma_{\vec \alpha_1}} & \text{by definition} \\
         \implies & \rr \arrow{\der_0}  \vec s s \ \notin\  \hr{ \sigma_{\vec \alpha_1}} & \text{for some $s \in \hr{\rho^{\vec \alpha_0}}$} \\
         \implies & \der_0 \vec s  s\  \notin \ \hr{ \sigma_{\vec \alpha_1}} & \text{by $\rr\arrow$ reduction} \\
         \implies & \der_0 \vec s s \notin \hr{ \sigma} & \text{by monotonicity}
    \end{array}
    \]
    where $\vec \alpha_0$ and $\vec \alpha_1$ are appropriate subsequences of $\vec \alpha$ in case not all fixed point subformulas of $\rho \to \sigma$ occur in $\rho$ or in $\sigma$.
    So we set $(\Sigma' \seqar \tau') := (\Sigma, \rho \seqar \sigma)$, and $\der' \dfn \der_0$ and $\vec s' := (\vec s,s)$.
    
    
    \item If $\rrule$ is
    \(
    \vliinf{\lr\arrow }{}{\Gamma,\Delta, \rho \to \sigma \seqar \tau}{\Gamma \seqar \rho}{\Delta , \sigma \seqar \tau}
    \)
    let $\vec r \in \hr{\Gamma}, \vec t\in \hr{\Delta}$ and $s \in \hr{\rho \to \sigma}$.
    Like before, let
    $\vec \alpha$ be the least assignment such that $s \in \hr{(\rho \to \sigma)^{\vec \alpha}}=\hr{\rho_{\vec \alpha_0}\to \sigma^{\vec \alpha_1}}$.
    We have two cases:
    \begin{itemize}
        \item if $\der_0\vec r \notin \hr{\rho_{\vec \alpha_0}}$ then we set $(\Sigma' \seqar \tau') \dfn (\Gamma \seqar \rho) $, and $\der'\dfn \der_0$ and $\vec s' \dfn \vec r$.
        \item otherwise ${\der_0} \vec r \in \hr{\rho_{\vec \alpha_0}}$ and we have:
        \[
        \begin{array}{rll}
             & {\lr \arrow}{\der_0}{\der_1} \vec r \vec t s \notin \hr{\tau} & \\
             \implies & {\der_1} \vec t (s ({\der_0} \vec r)) \notin \hr\tau & \text{by $\lr\arrow$ reduction}
        \end{array}
        \]
        Since $\der_0 \vec r \in \hr{\rho_{\vec \alpha_0}}$ and $s \in \hr{\rho_{\vec \alpha_0} \arrow \sigma^{\vec \alpha_1}}$ 
        we have $s( \der_0 \vec r) \in \hr{\sigma^{\vec\alpha_1}}$ so we set
        $(\Sigma' \seqar \tau') \dfn (\Delta, \sigma \seqar \tau)$, and $\der' \dfn \der_1$ and $\vec s' \dfn (\vec t, s(\der_0 \vec r) )$.
    \end{itemize}
    \item 
    If $\rrule$ is 
    \(
    \vlinf{\rr \mu}{}{\Sigma \seqar \mu X \sigma(X)}{\Sigma \seqar \sigma(\mu X \sigma(X))}
    \)
    and $\vec s \in \hr{\Sigma}$,
    \[
    \begin{array}{rll}
         &  \rr\mu\der_0 \vec s\  \notin \ \hr{\mu X\sigma(X)} \\
         \implies & \der_0 \vec s \notin \hr{\mu X\sigma(X)} & \text{by $\rr\mu$ reduction } \\
         \implies & \der_0 \vec s\  \notin\ \hr{\sigma(\mu X \sigma(X))} & \text{by Proposition~\ref{prop:hr-mu-is-lfp}}
    \end{array}
    \]
    so we set $(\Sigma' \seqar \tau') \dfn (\Sigma \seqar \sigma(\mu X \sigma (X)))$, $\der' \dfn \der_0$ and $\vec s' \dfn \vec s$.
    \item if $\rrule $ is 
    \(
    \vlinf{\mulunf}{}{\Gamma, \mu X \sigma(X) \seqar \tau}{\Gamma, \sigma(\mu X \sigma(X)) \seqar \tau}
    \)
    and $\vec r \in \hr{\Gamma}$ and $s \in \hr{\mu X \sigma(X)}$ we have:
    \[
    \begin{array}{rll}
         &  \lr \mu\der_0 \vec r s \ \notin \ \hr\tau \\
        \implies & \der_0 \vec r s \ \notin\  \hr \tau  & \text{by $\lr \mu$ reduction}
    \end{array}
    \]
    Since $s \in \hr{\mu X \sigma(X)}$ then also $s \in \hr{\sigma(\mu X \sigma(X))}$ by Proposition~\ref{prop:hr-mu-is-lfp},
    so we set $(\Sigma' \seqar \tau') \dfn (\Gamma , \sigma(\mu X \sigma(X)) \seqar \tau)$, and $ \der' \dfn \der_0$ and $\vec s' \dfn (\vec r , s)$.
    \end{itemize}
\end{definition}

   Note that the $\arrow$ cases for the Definition above are particularly subtle, requiring consideration of least ordinal assignments similar to the handling of $\lor$-left in fixed point modal logics, cf.\ e.g.\ \cite{NiwWal96:games-for-mu-calc}.


\anupam{make a remark about extension to $+$ cases and first component}


\subsection{Non-total branch construction}
From here the proof of the Interpretation Theorem \ref{thm:prog-implies-hr} proceeds by contradiction, as is usual in cyclic proof theory.
The definition above is used to construct an infinite `non-total' branch, along which there must be a progressing thread.
We assign ordinals approximating the critical fixed point formula and positive formulas of higher priority, which must always be present as positive subformulas on the LHS, or negative on the RHS, along the thread.
Tracking the definition above, we note that this ordinal assignment sequence is non-increasing; moreover at any $\mulunf$ step on the critical fixed point, the corresponding ordinal must be a successor and strictly decreases. 
Thus the ordinal sequence does not converge, yielding a contradiction.

% Further details are in Appendix~\ref{subsec:app-non-total-branch-construction}.




\anupam{commented remark about fixed points approximants below, pertinent to old proof, add similar remark back if time.}
% \begin{remark}
% Note that, in the above proof, it is not at all the case that each component $\alpha_{ij}$ is monotone non-increasing: it is possible that the fixed point formula $\mu X_j \sigma_j (X_j) $ `disappears' along the thread, and it may reappear later with a larger ordinal.
% However the definition of $<$ on types ensures that this is only when a fixed point of higher priority has been unfolded.
% In particular, since  $\mu X_1\sigma_1 (X_1)$ is a subformula of each $\rho_i$, for $i\geq k$, we do have that $(\alpha_{i1})_{i\geq k}$ is a monotone non-increasing sequence, which does not converge due to 
% \end{remark}


\begin{proof}
[Proof of \Cref{thm:prog-implies-hr}]
    Let $\der : \Sigma\seqar \tau$ and $\vec s \in \hr\Sigma$ be as in the statement, and suppose for contradiction that $\der\vec s \notin \hr{ \tau}.$
Setting $\der_0 \dfn \der$ and $\vec s_0 \dfn \vec s$, we use Definition~\ref{def:refl-non-totality} to construct an infinite branch $(\der_i : \Sigma_i \seqar \tau_i)_{i < \omega}$ and associated inputs $(\vec s_i)_{i < \omega}$ by always setting $\der_{i+1} \dfn \der_i'$, $(\Sigma_{i+1}\seqar \tau_{i+1}) \dfn (\Sigma_i'\seqar \tau_i')$ and $\vec s_{i+1}\dfn \vec s_i'$.
% let $\der_i$ and $\vec t_i$ be as constructed in the previous subsection,

Now let $(\rho_i)_{i\geq k}$ be a progressing thread along $(\der_i)_{i <\omega}$, since $\der$ is progressing, and let $r_i \in \vec s_i$ be the corresponding input, for $i\geq k$, when $\rho_i$ is on the LHS.
% 
% Without loss of generality assume that the positively occurring fixed points of each $\rho_i$ on the LHS and negatively occurring fixed points of each $\rho_i$ on the RHS are among $\mu X_1 \sigma_1(X_1) > \cdots > \mu X_n \sigma_n (X_n)$, for $i\geq k$,\footnote{This is because we have necessarily $\rho_{i+1}\flleq \rho_{i}$, and so the thread must eventually stabilise within some Fischer-Ladner class. Of course, no inference step changes polarity of the subformulas occurring on the same side.}
% 
Since $(\rho_i)_{i\geq k}$ is progressing, let $\mu X_1\sigma_1(X_1) > \cdots > \mu X_n \sigma_n (X_n)$ enumerate the fixed points occurring in every $\rho_i$ positively in the LHS, and negatively in the RHS, such that $\mu X_n\sigma_n$ is the smallest infinitely often principal formula along $(\rho_i)_{i\geq k}$. 
Note that such a finite set of fixed points must exist since 
necessarily $\rho_{i+1}\flleq \rho_{i}$, and so the thread must eventually stabilise within some Fischer-Ladner class. 
% Of course, no inference step changes polarity of the subformulas occurring on the same side.}
WLoG, no $\mu X_j \sigma_j (X_j)$, for $j<n$, is principal along $(\rho_i)_i$.
% Note that we have that $\mu X \sigma(X)$ is indeed a subformula of each $\rho_i$, positively on the LHS and negatively on RHS, and WLoG has highest priority among all fixed point formulas occurring.

By Proposition~\ref{prop:pos-neg-approximants}, let $\vec \alpha_i = \alpha_{i1}, \dots, \alpha_{in}$, for $i \geq k$, be the least assignments such that:
\begin{itemize}
    \item if $\rho_i$ is on the LHS then $r_i \in \hr{\rho_i^{\vec \alpha_i}}$; 
    \item if $\rho_i$ is on the RHS then $\der_i\vec s_i \notin (\rho_i)_{\vec \alpha_i}$. 
\end{itemize}

We claim that $(\vec \alpha_i)_{i\geq k}$ is a monotone non-increasing sequence of ordinal assignments that does not converge:
\begin{itemize}
    \item By construction of $\der_i$ and $\vec s_i$, appealing to Definition~\ref{def:refl-non-totality}, note that for each step for which $\mu X_n\sigma_n(X_n)$ is not principal (along $(\rho_i)_i$ on the LHS) we have that $\vec \alpha_{i+1}\leq \vec \alpha_i$.
    Note in particular that the $\arrow$-cases of Definition~\ref{def:refl-non-totality} are designed to guarantee $\vec \alpha_{i+1}\leq \vec \alpha_i$ at $\arrow$-steps.
    \item Now, consider a $\mulunf$-step along the progressing thread on the critical fixed point formula,
    \[
    \vlinf{\mulunf}{}{\Gamma, \mu X_n \sigma_n(X_n)\seqar \pi}{\Gamma , \sigma_n(\mu X_n \sigma_n(X_n))\seqar \pi}
    \]
    with $\rho_i = \mu X_n \sigma_n (X_n)$.
    Writing $\vec \alpha_{in} = \alpha_{i,1}, \dots, \alpha_{i,n-1}$ (a prefix of $\vec \alpha_i$) we have:
    \[
    \begin{array}{rl}
        \hr{(\mu X_n \sigma_n (X_n))^{\vec \alpha_{in}\alpha_{in}}} 
        = & \hr{(\mu X_n \sigma_n^{\vec \alpha_{in}} (X_n))^{\alpha_{in}}} \\
        = & \hr{\sigma_n^{\vec \alpha_{in} \alpha_{in}}(\emptyset)} \\
        = & \hr{\sigma_n^{\vec \alpha_{in} \succ \alpha_{in}'}(\emptyset)} \\
        = & \hr{\sigma_n^{\vec \alpha_{in}} (\hr{\sigma_n^{\vec \alpha_{in} \alpha_{in}' } (\emptyset) } ) } \\
         = & \hr{\sigma_n^{\vec \alpha_{in}} (\hr{(\mu X_n \sigma_n^{\vec \alpha_{in}} (X_n))^{\alpha_{in}'}} ) } \\
           = & \hr{\sigma_n^{\vec \alpha_{in}} ( \hr{(\mu X_n \sigma_n(X_n))^{\vec \alpha_{in}\alpha_{in}'}}) } \\
    \end{array}
    \]
    % This means that $t \in \hr{(\mu X_j \sigma_j (X_j))^{}} \implies t\in \hr{}$
    for some $\alpha_{in}'$ since $\alpha_{ij}$ must be a successor ordinal.
    % \footnote{An element must enter a fixed point for the first time at a successor ordinal, since a limit approximant is just the union of smaller approximants.}
    Thus indeed $\vec \alpha_{i+1} < \vec \alpha_{i}$.
\end{itemize}
% 
% Let $(\alpha_i)_{i\geq k}$ be the least ordinals such that $t_i \in \hr{\sigma_i^{\alpha_i}}$.
% We claim that $(\alpha_i)_{i\geq k}$ is a monotone non-increasing sequence that does not converge.\oldtodo{at this point can make do with only a single ordinal, only need ordinal sequences for the definition of the branch.}
% 
% First, note that for $i \geq k$, each $\sigma_i$ has the same FL-closure, and so assignments approximate the same fixed-point subformulas, say $\mu X_1 \sigma_1(X_1) > \cdots > \mu X \sigma (X)$
% 
% First write $\alpha_i = \vec \alpha_i$ (construed as a sequence).
% Note that each $\vec \alpha_i$ and $\vec \alpha_j$ have same length for $i,j\geq k$, since such $\sigma_i, \sigma_j$ must be in the same Fischer-Ladner class.\oldtodo{justify this WLoG}
% \begin{itemize}
%     \item By the construction of $D_i$ and $\vec t_i$, we have that always $\alpha_{i+1}\leq \alpha_i$.\oldtodo{case analysis}
%     \item At a left $\mu$ step,\oldtodo{again, turn this into lemma so that notation/subscripts do not clash}
%     \[
%     \vlinf{\mu l}{}{\vec \sigma, \mu X \sigma(X)\seqar \tau}{\vec \sigma , \sigma(\mu X \sigma(X))\seqar \tau}
%     \]
%     suppose we have $\vec t \in \hr{\vec \sigma^{\vec \alpha}}$, $t\in \hr{(\mu X \sigma(X))^\alpha}$ such that $\sem{D}\vec t t \notin \hr{\tau}$.
%     We claim that there is $\alpha'<\alpha$ such that $t\in \hr{\sigma(\mu X \sigma(X))^{\alpha'}}$.
    
%     First, let $\alpha = (\vec \beta, \gamma, \vec \delta)$ such that $\gamma$ corresponds to the fixed point $\mu X \sigma(X)$.
% \end{itemize}
% \oldtodo{finish}
\noindent
This contradicts the well-foundedness of ordinals.
\end{proof}